\documentclass[10pt,a4paper]{article}

\usepackage[utf8]{inputenc}
\usepackage[T1]{fontenc}
\usepackage[spanish,es-nodecimaldot]{babel}
\usepackage{tgheros}
\renewcommand{\familydefault}{\sfdefault}
\usepackage[a4paper,top=1.85cm,bottom=1.8cm,left=1.9cm,right=1.9cm,headheight=20pt]{geometry}
\usepackage{xcolor}
\usepackage{tabularx,array,booktabs,colortbl}
\usepackage{enumitem}
\usepackage[most]{tcolorbox}
\usepackage{fancyhdr}
\usepackage{lastpage}
\usepackage{microtype}
\usepackage{setspace}
\usepackage{titlesec}
\usepackage{hyperref}
\usepackage{xurl}
\usepackage{tikz}
\usepackage{ragged2e}

% ---------- PALETA ROUTTECH ----------
\definecolor{RTForest}{HTML}{25463B}
\definecolor{RTTerracotta}{HTML}{B45F45}
\definecolor{RTGold}{HTML}{D5A64A}
\definecolor{RTCream}{HTML}{F7F2E8}
\definecolor{RTLight}{HTML}{EEF3F0}
\definecolor{RTGray}{HTML}{4D5652}

% ---------- DATOS ----------
\newcommand{\Proyecto}{RoutTech: Sistema de Movilidad Colaborativa para la Comunidad Universitaria de la UTEQ}
\newcommand{\Asignatura}{ISR-401 -- Ingeniería de Requisitos}
\newcommand{\Periodo}{2026--2027 PPA}
\newcommand{\CategoriaEtica}{Categoría B -- Datos personales}
\newcommand{\Repositorio}{https://github.com/AnderJM3/RoutTech_ISR401}
\newcommand{\FechaDocumento}{29 de julio de 2026}
\newcommand{\VersionDocumento}{2.0}
\newcommand{\Paralelo}{Cuarto nivel -- Software B}
\newcommand{\EstadoDocumento}{Versión 2.0 para firma}
\newcommand{\Mayra}{CORDOVA CARREÑO MAYRA LUCILA}
\newcommand{\Anderson}{NARANJO FLORES ANDERSON JEAMPIERE}
\newcommand{\Docente}{Ing. Gleiston Guerrero Ulloa, PhD}

% ---------- TABLAS ----------
\newcolumntype{Y}{>{\RaggedRight\arraybackslash}X}
\newcolumntype{C}[1]{>{\centering\arraybackslash}p{#1}}
\newcolumntype{L}[1]{>{\RaggedRight\arraybackslash}p{#1}}
\renewcommand{\arraystretch}{1.22}
\setlength{\tabcolsep}{5pt}

% ---------- ENCABEZADO Y PIE ----------
\pagestyle{fancy}
\fancyhf{}
\fancyhead[L]{\small\textbf{UTEQ | FCC | ISR-401}}
\fancyhead[R]{\small\textbf{B3 -- Datos sintéticos}}
\fancyfoot[L]{\small Paquete ético RoutTech}
\fancyfoot[R]{\small Página \thepage\ de \pageref{LastPage}}
\renewcommand{\headrulewidth}{0.8pt}
\renewcommand{\footrulewidth}{0.4pt}
\renewcommand{\headrule}{\hbox to\headwidth{\color{RTTerracotta}\leaders\hrule height \headrulewidth\hfill}}
\renewcommand{\footrule}{\hbox to\headwidth{\color{RTGold}\leaders\hrule height \footrulewidth\hfill}}

% ---------- TÍTULOS ----------
\titleformat{\section}
  {\large\bfseries\color{RTForest}}
  {\thesection.}{0.45em}{}
  [\vspace{1pt}\color{RTGold}\titlerule]
\titlespacing*{\section}{0pt}{0.7em}{0.35em}

\hypersetup{
  colorlinks=true,
  linkcolor=RTForest,
  urlcolor=RTTerracotta,
  pdftitle={B3 Compromiso de Datos Sinteticos RoutTech},
  pdfauthor={Equipo RoutTech}
}
\urlstyle{same}
\setlength{\parindent}{1.1em}
\setlength{\parskip}{2.5pt}
\setstretch{1.10}
\setlist[enumerate]{leftmargin=1.6em,itemsep=3pt,topsep=3pt}

\begin{document}
\justifying

% ======================= PORTADA =======================
\thispagestyle{empty}
\begin{tikzpicture}[remember picture,overlay]
  \fill[RTForest] (current page.north west) rectangle ([yshift=-4.6cm]current page.north east);
  \fill[RTTerracotta] ([yshift=-4.6cm]current page.north west) rectangle ([yshift=-4.95cm]current page.north east);
\end{tikzpicture}

\vspace*{0.65cm}
\begin{center}
{\color{white}\Large\bfseries UNIVERSIDAD TÉCNICA ESTATAL DE QUEVEDO}

\vspace{1.75cm}
{\large\bfseries\color{RTForest} Facultad de Ciencias de la Computación\\[4pt]}
{\normalsize\color{RTGray} Carrera de Ingeniería de Software}

\vspace{0.8cm}

{\Huge\bfseries\color{RTForest} COMPROMISO DE DATOS SINTÉTICOS\\[8pt]}
{\Large\bfseries\color{RTTerracotta} B3 -- Categoría B}

\vspace{1.0cm}

\begin{tcolorbox}[width=0.90\textwidth,colback=RTCream,colframe=RTGold,boxrule=1pt,arc=2mm]
\centering
{\large\bfseries \Proyecto}\\[7pt]
\textbf{Clasificación ética:} \CategoriaEtica\\
\textbf{Asignatura:} \Asignatura\\
\textbf{Periodo académico:} \Periodo
\end{tcolorbox}

\vfill
\begin{tabular}{rl}
\textbf{Documento:} & B3 -- Compromiso de uso exclusivo de datos sintéticos\\
\textbf{Versión:} & \VersionDocumento\\
\textbf{Estado:} & \EstadoDocumento\\
\textbf{Paralelo:} & \Paralelo\\
\textbf{Repositorio:} & \href{\Repositorio}{RoutTech\_ISR401}\\
\textbf{Fecha:} & \FechaDocumento
\end{tabular}

\vspace{0.75cm}
{\small Quevedo -- Los Ríos -- Ecuador}
\end{center}
\clearpage

% ======================= CONTENIDO =======================
\section{Identificación}

\rowcolors{2}{RTLight}{white}
\noindent\begin{tabularx}{\textwidth}{L{4.5cm} Y}
\rowcolor{RTForest}
\textcolor{white}{\textbf{Campo}} & \textcolor{white}{\textbf{Información}}\\
Proyecto & \Proyecto.\\
Equipo investigador & \Mayra\ y \Anderson.\\
Docente responsable & \Docente.\\
Categoría ética & \CategoriaEtica.\\
Versión y fecha & Versión \VersionDocumento, \FechaDocumento.\\
\end{tabularx}

\section{Declaración de compromiso}

Los integrantes del proyecto RoutTech declaramos que, durante el desarrollo académico, la construcción del MVP, las pruebas, las demostraciones y la elaboración de evidencias se realizarán exclusivamente con datos ficticios, sintéticos, agregados o previamente disociados. En consecuencia, asumimos los siguientes compromisos:

\begin{enumerate}
  \item No importar ni copiar bases reales de viajes, reservas, pagos, credenciales institucionales, vehículos, placas, calificaciones, historiales ni ubicaciones individuales.
  \item Utilizar personas ficticias, correos de prueba, placas inventadas, rutas generales, horarios simulados y reservas sintéticas en el MVP, capturas, demostraciones y casos de prueba.
  \item No emplear nombres, cédulas, teléfonos, correos personales, domicilios, coordenadas exactas ni trayectorias reales de participantes.
  \item Analizar las respuestas del cuestionario de forma agregada y mantener las entrevistas seudonimizadas mediante códigos de evidencia.
  \item Usar únicamente información previamente disociada cuando una unidad institucional facilite datos para fines académicos.
  \item Aislar cualquier dato real capturado por error, comunicar el incidente al docente responsable y eliminarlo de forma segura dentro de las 24 horas siguientes a su detección.
  \item No implementar pagos reales, seguimiento de ubicación en tiempo real, autenticación institucional productiva ni conexión con sistemas oficiales de la UTEQ durante esta fase.
  \item Revisar antes de publicar que los documentos, bases de prueba, capturas y demostraciones no permitan identificar directa o indirectamente a una persona.
\end{enumerate}

\section{Datos excluidos y datos permitidos}

\rowcolors{2}{RTCream}{white}
\noindent\begin{tabularx}{\textwidth}{Y Y}
\rowcolor{RTTerracotta}
\textcolor{white}{\textbf{No se utilizarán datos reales de}} &
\textcolor{white}{\textbf{Se utilizarán únicamente}}\\
Personas, credenciales, placas, vehículos, reservas, pagos, calificaciones, ubicaciones exactas e historiales individuales. &
Personas ficticias, códigos de prueba, placas inventadas, zonas generales, rutas simuladas, reservas sintéticas y resultados agregados.\\
\end{tabularx}

\section{Firmas}

\vspace{0.35cm}
\noindent\begin{tabularx}{\textwidth}{Y C{4.4cm} C{3.1cm}}
\rowcolor{RTForest}
\textcolor{white}{\textbf{Nombre}} &
\textcolor{white}{\textbf{Firma}} &
\textcolor{white}{\textbf{Fecha}}\\
\Mayra & \rule{4.0cm}{0.4pt} & \rule{2.6cm}{0.4pt}\\[5pt]
\Anderson & \rule{4.0cm}{0.4pt} & \rule{2.6cm}{0.4pt}\\[5pt]
\Docente & \rule{4.0cm}{0.4pt} & \rule{2.6cm}{0.4pt}\\
\end{tabularx}

\end{document}
